\section{DSYSTEM to FO-LTL Encoding}
\label{app:encoding}

This appendix presents the encoding of Section~\ref{sec:dsystem-encoding}. The bi-implication below is also machine-checked in a companion formalization.

Given $Z=(S,I,O,\delta,\rho)$ and $s_0\in S$,
\[
  \texttt{compileSystemFO}(Z,s_0) \;:=\; \text{init}(s_0)\wedge\text{step}(Z)\wedge\text{readout}(Z)
\]
with the constructors of \eqref{eq:fo-init}--\eqref{eq:fo-readout}.

\begin{quote}
\textbf{Theorem (execution equivalence).} Fix $Z$, $s_0$, and $f$. Trajectories $(g,y)$ form a valid Wymore execution from $s_0$ under $f$ if and only if $(f,g,y)\models_{\mathrm{FO}}\texttt{compileSystemFO}(Z,s_0)$.
\end{quote}

\paragraph{Proof ($\Rightarrow$).} Canonical trajectories satisfy init, step, and readout by Wymore's construction.

\paragraph{Proof ($\Leftarrow$).} Satisfaction of init and step, with trajectory uniqueness, identifies $g$ with the canonical state trajectory; readout identifies $y$.

When $|S|$, $|I|$, and $|O|$ are finite, a propositional LTL encoding of the transition table exists as a syntactic corollary; it is an expressibility result, not a practical encoder for large products $|S|\cdot|I|$.
